\documentclass[print,thumb,showtrims]{personal}
% options:
%   british (default), american -- language: hyphenation, quotation marks, dates
%   a5 (default), royal, tudelft -- book page size: 148x210, 156x234 or 170x240 mm
%   print     -- print as a book
%   thumb     -- print thumb index, only available for print
%   showtrims -- show trim section, only available for print
%   showframe -- show construction guide (page layout)

% packages

%\usepackage[utf8]{inputenc} % already in the cls file
% \usepackage{graphicx}   % https://ctan.org/pkg/graphicx %% provides optional arguments to the \includegraphics command
% \usepackage{tikz}       % https://www.ctan.org/pkg/pgf %% TikZ manual found here, allows image drawing
\usepackage{bm}         % https://ctan.org/pkg/bm %% more powerful package to provide bold symbol support (use \bm{} instead of \boldsymbol{})
\usepackage[version=4]{mhchem} % https://ctan.org/pkg/mhchem?lang=en %% gives ability to implement chemical formulas via the \ce command
\usepackage{siunitx}    % https://ctan.org/pkg/siunitx %% Provides SI unit control through the \SI{<value>}{<units>} command
%\usepackage{longtable}  % https://ctan.org/pkg/longtable %% gives the longtable environment, which allows multiple pages!
%\usepackage{tabularx}   % https://ctan.org/pkg/tabularx %% gives the tabularx environment, which allows multiple lines within a cell without super weird coding
\usepackage{xltabular}  % https://ctan.org/pkg/xltabular %% Gives the functionality of both longtable + tabularx package together, check https://www.overleaf.com/latex/examples/a-longtable-plus-tabularx-example/rcrxtjbrvyft for an example of how it works.
% \usepackage{hyperref} % if you want to add external links to words, check userguide for info % already in the cls file
\usepackage{float}      % https://ctan.org/pkg/float %% Allows figure forcing spots ([H], etc.)
\usepackage{multicol}   % https://ctan.org/pkg/multicol %% useful for multiple column pages. \multicols gives a balanced output between each column, whereas \multicols* gives unbalanced.
\setlength{\columnsep}{1cm}
\usepackage{appendix}   % https://ctan.org/pkg/appendix %% The appendix package provides various ways of formatting the titles of appendices. 
%\usepackage{caption}    % https://ctan.org/pkg/caption %% allows extensive control of captions. Does not come with subcaption (see Chap 6 of manual), already in cls file
\usepackage{subcaption} % https://ctan.org/pkg/subcaption %% provides subfigure/caption support i.e use: \begin{subfigure}{\textwidth}
\usepackage{array}      % https://ctan.org/pkg/array %% gives more power to making tables! array provides many more utilities to making tables
\usepackage{multirow}   % https://ctan.org/pkg/multirow %% centering of text over multiple rows for tables!
\usepackage{mathtools}  % https://ctan.org/pkg/mathtools %% allows the use of \splitfrac{}{}, splitdfrac{}{}
%\usepackage[skip]{parskip}   % https://ctan.org/pkg/parskip %% Changes from indentations to line skips between paragraphs % already included in the memoir class, use lines below instead
\nonzeroparskip
\setlength{\parindent}{0pt}
% \usepackage{chngcntr}   % https://ctan.org/pkg/chngcntr %% Change the resetting of counters by using \counterwithin (which sets up a counter to be reset when another is incremented) and \counterwithout (which unsets such a relationship). %%% ALREADY IN MEMOIR CLASS
% \usepackage{wrapfig}    % https://ctan.org/pkg/wrapfig %% allows the use of 'wrapfigure' environments
\usepackage{pdfpages}   % https://ctan.org/pkg/pdfpages %% This package simplifies the inclusion of external multi-page PDF documents in LaTeX documents. 
% \usepackage{xcolor}     % https://ctan.org/pkg/xcolor %% Provides easy driver-independent access to several kinds of color tints, shades, tones, and mixes of arbitrary colors.
\usepackage{colortbl}   % https://www.ctan.org/pkg/colortbl %% The package allows rows and columns to be coloured, and even individual cells.
\usepackage{booktabs}   % https://ctan.org/pkg/booktabs/ %% The package enhances the quality of tables in LaTeX
% \usepackage{arydshln}   % https://www.ctan.org/pkg/arydshln %% The package is to draw dash-lines in array/tabular environments.
\usepackage{listings}   % https://www.ctan.org/pkg/listings %% Use this for code
\usepackage{bigints}    % https://www.ctan.org/pkg/bigints %% \bigintssss, \bigintsss, \bigintss, \bigints, and \bigint
\usepackage{cancel}     % https://www.ctan.org/pkg/cancel %% A package to draw diagonal lines (“cancelling” a term) and arrows with limits (cancelling a term “to a value”) through parts of maths formulae
\usepackage{stackengine}% https://www.ctan.org/tex-archive/macros/latex/contrib/stackengine %% The stackengine package provides a versatile way to vertically stack objects in a variety of customizable ways.
%\usepackage{arydshln}   % https://ctan.org/pkg/arydshln %% The package is to draw dash-lines in array/tabular environments. Horizontal lines are drawn by \hdashline and \cdashline while vertical ones can be specified as a part of the preamble using ‘:’. The shape of dash-lines may be controlled through style parameters or optional arguments. %% CONFLICT WITH BOOKTABS
\usepackage{pifont}     % https://www.ctan.org/pkg/pifont %% The package provides commands for Pi fonts (Dingbats, Symbol, etc.); all commands assume you know the character number within the font of the symbol within the font. 
\usepackage{ragged2e}   % https://www.ctan.org/pkg/ragged2e %% The package defines new commands \Centering, \RaggedLeft, and \RaggedRight and new environments Center, FlushLeft, and FlushRight, which set ragged text and are easily configurable to allow hyphenation (the corresponding commands in LaTeX, all of whose names are lower-case, prevent hyphenation altogether)
\usepackage{bookmark}   % https://ctan.org/pkg/bookmark?lang=en %% This package implements a new bookmark (outline) organization for package hyperref. Bookmark properties such as style and color can now be set. Other action types are available (URI, GoToR, Named). The bookmarks are generated in the first compile run. Package hyperref uses two runs. 
\usepackage{nicematrix} %https://ctan.org/pkg/nicematrix %% This package is based on the package array. It creates PGF/TikZ nodes under the cells of the array and uses these nodes to provide functionalities to construct tabulars, arrays and matrices. 
\usepackage{notoccite}  % https://www.ctan.org/pkg/notoccite %% If you have \cite commands in \section-like commands, or in \caption, the citation will also appear in the table of contents, or list of whatever. If you are also using an unsrt-like bibliography style, these citations will come at the very start of the bibliography, which is confusing. This package suppresses the effect.
\usepackage{pdflscape}  % https://www.ctan.org/pkg/pdflscape %% The package adds PDF support to the landscape environment of package lscape, by setting the PDF /Rotate page attribute. Pages with this attribute will be displayed in landscape orientation by conforming PDF viewers.
% package fixes
\makeatletter % fix provided by https://tex.stackexchange.com/questions/522920/xltabular-breaking-booktabs
\def\@BTrule[#1]{%
  \ifx\longtable\undefined
    \let\@BTswitch\@BTnormal
  \else\ifx\hline\LT@hline
    \nobreak
    \let\@BTswitch\@BLTrule
  \else
     \let\@BTswitch\@BTnormal
  \fi\fi
  \global\@thisrulewidth=#1\relax
  \ifnum\@thisruleclass=\tw@\vskip\@aboverulesep\else
  \ifnum\@lastruleclass=\z@\vskip\@aboverulesep\else
  \ifnum\@lastruleclass=\@ne\vskip\doublerulesep\fi\fi\fi
  \@BTswitch}
\makeatother 
% code settings
\definecolor{dkgreen}{rgb}{0,0.6,0}
\definecolor{gray}{rgb}{0.5,0.5,0.5}
\definecolor{mauve}{rgb}{0.58,0,0.82}
% https://tex.stackexchange.com/questions/24528/having-problems-with-listings-and-utf-8-can-it-be-fixed
\lstset{
    frame=tb,
    language=Python,
    aboveskip=3mm,
    belowskip=3mm,
    showstringspaces=false,
    columns=flexible,
    basicstyle={\small\ttfamily},
    numbers=left,
    numbersep=5pt,
    numberstyle=\tiny\color{gray},
    keywordstyle=\color{blue},
    commentstyle=\color{dkgreen},
    stringstyle=\color{mauve},
    breaklines=true,
    breakatwhitespace=false,
    tabsize=3,
    inputencoding=utf8,
    extendedchars=true,
    literate={á}{{\'a}}1 {ã}{{\~a}}1 {é}{{\'e}}1 {ü}{{\"u}}1,
}
% page styles
%% Page styles (memoir). Running heads in letterspaced capitals, page numbers in the same size
%% (\runningheadtext and \runningheadfolio in fontText.tex). New page styles have empty heads and feet.
%%   mainmatter     : chapter title on left pages, section on right pages, thumb tabs
%%   chapteronly    : chapter title on both pages (chapters without sections, the references)
%%   appendix       : "Appendix A - title" on both pages, thumb tabs
%%   chapteropening : chapter opening pages: only the thumb tab
%%   frontmatter    : page number centred at the bottom
\makepagestyle{mainmatter}
\makeevenhead{mainmatter}{\runningheadfolio\lthumb\thepage}{\runningheadtext{\leftmark}}{}
\makeoddhead{mainmatter}{}{\runningheadtext{\rightmark}}{\runningheadfolio\thepage\rthumb}
\makeheadrule{mainmatter}{\textwidth}{0.3pt}

\makepagestyle{appendix}
\makeevenhead{appendix}{\runningheadfolio\lthumb\thepage}{\runningheadtext{Appendix \thechapter~\textendash~\leftmark}}{}
\makeoddhead{appendix}{}{\runningheadtext{Appendix \thechapter~\textendash~\leftmark}}{\runningheadfolio\thepage\rthumb}
\makeheadrule{appendix}{\textwidth}{0.3pt}

%% Chapter opening pages: no running head, but the thumb tab (the page a reader is looking for).
\makepagestyle{chapteropening}
\makeevenhead{chapteropening}{\lthumb}{}{}
\makeoddhead{chapteropening}{}{}{\rthumb}

\makepagestyle{frontmatter}
\makeevenhead{frontmatter}{}{}{}
\makeoddhead{frontmatter}{}{}{}
\makeevenfoot{frontmatter}{}{\runningheadfolio\thepage}{}
\makeoddfoot{frontmatter}{}{\runningheadfolio\thepage}{}

%% plain: used for chapter opening pages in the front and back matter (page number at the bottom)
\makeevenfoot{plain}{}{\runningheadfolio\thepage}{}
\makeoddfoot{plain}{}{\runningheadfolio\thepage}{}

\makepagestyle{chapteronly}
\makeevenhead{chapteronly}{\runningheadfolio\lthumb\thepage}{\runningheadtext{\leftmark}}{}
\makeoddhead{chapteronly}{}{\runningheadtext{\leftmark}}{\runningheadfolio\thepage\rthumb}
\makeheadrule{chapteronly}{\textwidth}{0.3pt}

\renewcommand{\chaptermark}[1]{\markboth{#1}{}}
%% The section number is stored together with its title, so the running head never pairs
%% the title of one section with the number of another.
\renewcommand{\sectionmark}[1]{\markright{\thesection.\enspace#1}}

% new commands
\newcommand{\NoLink}[1]{{\protect\NoHyper#1\protect\endNoHyper}} % Allows the use of \citeauthor without actually referencing back to the reference list, in case you want to just say the author's name. Example: \NoLink{\citeauthor{CoolKid2019}}
\newcommand\invisiblechapter[1]{                % makes an invisible chapter!
  \refstepcounter{chapter}%
  \addcontentsline{toc}{chapter}{\protect\numberline{\thechapter}#1}%
  \chaptermark{#1}}
\newcommand\invisiblesection[1]{                % makes an invisible section!
  \refstepcounter{section}%
  \addcontentsline{toc}{section}{\protect\numberline{\thesection}#1}%
  \sectionmark{#1}}
\newcommand\invisiblesubsection[1]{%            % makes an invisible subsection!
  \refstepcounter{subsection}%
  \addcontentsline{toc}{subsection}{\protect\numberline{\thesubsection}#1}%
  \subsectionmark{#1}}

\newcommand\diff[2][]{\mathrm{d}^{#1}{#2}}      % creates a nth-order ordinary differential variable
\newcommand\pdiff[2][]{\partial^{#1}{#2}}       % creates a nth-order partial differential variable
\newcommand\sdiff[2][]{\mathrm{D}^{#1}{#2}}     % creates a nth-order substantial differential variable
\newcommand\deriv[3][]{                         % creates a nth-order ordinary derivative
\frac{\diff[#1]{#2}}{\diff{#3}^{#1}}
}
\newcommand\pderiv[3][]{                        % creates a nth-order partial derivative
\frac{\pdiff[#1]{#2}}{\pdiff{#3}^{#1}}
}
\newcommand\sderiv[3][]{                        % creates a nth-order substantial derivative
\frac{\sdiff[#1]{#2}}{\sdiff{#3}^{#1}}
}
\newcommand\cderiv[3][]{\diff[#1]{}_{#3}{#2}}   % creates a compact form nth-order ordinary derivative
\newcommand\cpderiv[3][]{\pdiff[#1]{}_{#3}{#2}} % creates a compact form nth-order partial derivative
\newcommand\csderiv[3][]{\sdiff[#1]{}_{#3}{#2}} % creates a compact form nth-order substantial derivative

\newcommand\uVec[2]{\hat{\Vec{#1}}_{#2}}        % creates a unit vector
\newcommand\uVecT[2]{\overline{\Vec{#1}}_{#2}}  % creates a transformed unit vector
\newcommand\transpose{^\textsf{T}}              % places in a transpose symbol

%https://tex.stackexchange.com/questions/417310/double-bar-notation-for-second-order-tensors4
\newcommand\barbelow[1]{\stackunder[1.6pt]{$#1$}% creates a bar below an object
                        {\rule{0.9ex}{.1ex}}}   

\newcommand{\grad}{ \barbelow{\nabla}}          % \grad now includes bar below.
\newcommand{\curl}{ \barbelow{\nabla}\times}    % \curl now includes bar below.
\newcommand\proj[2]{ \mathsf{proj}_{#1}^{#2}\,} % Projection operator
\newcommand{\Res}{ \mathsf{Res}}                % Residual operator

%% Create typical dimensionless numbers used in fluid + heat transfer problems
\renewcommand{\Pr}{\mathrm{Pr}}                 % Prandtl number
\renewcommand{\Re}{\mathrm{Re}}                 % Reynolds number
\newcommand{\Fr}{\mathrm{Fr}}                   % Froude number
\newcommand{\Eu}{\mathrm{Eu}}                   % Euler number
\newcommand{\St}{\mathrm{St}}                   % Strouhal number
\newcommand{\Gr}{\mathrm{Gr}}                   % Grashof number
\newcommand{\Ma}{\mathrm{Ma}}                   % Mach number
\newcommand{\Pe}{\mathrm{Pe}}                   % Péclet number
\newcommand{\Tu}{\mathrm{Tu}}                   % Turbulence level

%% With Eulervm, no \therefore defined. newtxmath provides \therefore, so this is only defined if missing.
% https://tex.stackexchange.com/questions/142551/how-to-put-therefore-and-implies-symbols
\providecommand{\therefore}{\dot{.\hspace{.055in}.}\hspace{.2in}}

%% standard /mapsto symbol, but reversed
% https://tex.stackexchange.com/questions/26508/left-version-of-mapsto
\providecommand{\mapsfrom}{\mathrel{\reflectbox{\ensuremath{\mapsto}}}}
\providecommand{\longmapsfrom}{\mathrel{\reflectbox{\ensuremath{\longmapsto}}}}

%% Using the same figure twice with no new number
% https://tex.stackexchange.com/questions/200211/using-the-same-figure-twice-with-no-new-number
\newcommand{\repeatcaption}[2]{%
  \renewcommand{\thefigure}{\ref{#1}}%
  \captionsetup{list=no}%
  \caption{#2 (repeated from page \pageref{#1})}%
  \addtocounter{figure}{-1}% So that next figure after the repeat gets the right number.
}

% Create delimiter environments (e.g floor(t) produces ⌊t⌋)
% --------------------
% https://tex.stackexchange.com/questions/118173/how-to-write-ceil-and-floor-in-latex
\DeclarePairedDelimiter\ceil{\lceil}{\rceil}    % Use ceil*{...} to produce big brackets
\DeclarePairedDelimiter\floor{\lfloor}{\rfloor} % Use floor*{...} to produce big brackets

% Create new math operators
% --------------------
\DeclareMathOperator*{\argmax}{arg\,max}
\DeclareMathOperator*{\argmin}{arg\,min}

% tensor - double dot product
\newcommand{\cddot}{\mathbin{{\cdot}{\cdot}}}

% Renew Commands:
% --------------------
\renewcommand{\vec}[1]{ \barbelow{\bm{#1}} }    % \vec to \bm combination
\renewcommand{\div}{ \barbelow{\nabla}\cdot}    % \div now includes bar below.

\renewcommand*{\tocheadstart}{\vspace*{0pt}}    % straight from the memoir manual, may need to change if you change
\renewcommand*{\lofheadstart}{\vspace*{0pt}}    % titlespacing for chapters... investigate?
\renewcommand*{\lotheadstart}{\vspace*{0pt}}    % pre-title spacing is now the same for toc, lof, lot and nomenclature

% \itemize labels %% https://texblog.org/2008/10/16/lists-enumerate-itemize-description-and-how-to-change-them/
\renewcommand{\labelitemi}{-}
\renewcommand{\labelitemii}{-}
\renewcommand{\labelitemiii}{-}
\renewcommand{\labelitemiv}{-}

% Robust Commands:
% --------------------
\DeclareRobustCommand{\VAN}[0]{Van} % Within text, authors with a van/Van prefix use Van within text, van outside text. 


%% Enable the following to display watermarks
% --------------------
% \usepackage{draftwatermark}   % https://www.ctan.org/pkg/draftwatermark %% Create watermarks for your file
% \SetWatermarkLightness{0.8}
% \SetWatermarkText{DRAFT}
% --------------------

%% Use the following packages to provide lipsum paragraphs or visualize linenumbers.
\usepackage{lipsum} % disable when not needed
% \usepackage{lineno}
% \linenumbers


% Useful Links:
% https://www.sciencedirect.com/science/article/pii/B9780080982854000042 - How to write a literature review
% http://web.mit.edu/rsi/www/pdfs/advmath.pdf - LaTeX Math Tools
% https://en.wikibooks.org/wiki/LaTeX/Advanced_Mathematics - More LaTeX Math Tools

\AtBeginDocument{\colorlet{tudelft-black}{.}}
\begin{document}
% \sisetup{   detect-family,
%             detect-weight,
%             math-rm = \mathrm }

% Force recalculation of the current page node whenever the overlay option is used.
% If your page is centered on your stock paper it is only necessary to calculate
% the current page rectangle once.
% https://texample.net/tikz/examples/modifying-current-page-node/
\pgfkeys{/tikz/overlay/.add code={}{\setpagenode}}

\pagenumbering{gobble}
% \thispagestyle{empty}
\vspace*{\fill}
\begin{center}
    {\LARGE \textit{This page is intentionally left blank.}}
\end{center}
\vspace*{\fill}
\clearpage

%% Title pages of a TU Delft dissertation. All names, dates and numbers come from _thesisinfo.tex.
%%   page i   : half-title
%%   page ii  : blank
%%   page iii : official title page (wording prescribed by TU Delft; check it against the
%%              Graduate School's current template before printing)
%%   page iv  : doctoral committee and colophon

%% ---------------------------------------------------------------------
%% Page i: half-title (title only, smaller than on the title page, at the height of the chapter titles)
%% ---------------------------------------------------------------------
\chapterstyle{frontmatter} % the chapter-title spacing, so the titles sit at the height of the chapter titles
\thispagestyle{empty}
\makeatletter
\begingroup
    \parskip\z@ \parindent\z@
    \chapterheadstart\thesis@chapterlabelspace
    \centering
    {\fontsize{0.75\dimexpr\chaptertitlesize\relax}{1.15\dimexpr\chaptertitlesize\relax}\headingfont\thesistitle\par}
    \medskip
    {\headingfont\large\thesissubtitle\par}
\endgroup
\makeatother

%% ---------------------------------------------------------------------
%% Page ii: blank
%% ---------------------------------------------------------------------
\cleardoublepage

%% ---------------------------------------------------------------------
%% Page iii: official title page
%% ---------------------------------------------------------------------
\thispagestyle{empty}
\makeatletter
\begingroup
    \parskip\z@ \parindent\z@
    \chapterheadstart\thesis@chapterlabelspace
    \centering
    {\fontsize{\chaptertitlesize}{1.15\dimexpr\chaptertitlesize\relax}\headingfont\thesistitle\par}
    \bigskip
    {\headingfont\Large\thesissubtitle\par}
\endgroup
\vfill
\begin{center}
    {\runningheadcaps DISSERTATION}

    \bigskip

    for the purpose of obtaining the degree of doctor \\
    at Delft University of Technology \\
    by the authority of the Rector Magnificus, \rectormagnificus, \\
    chair of the Board for Doctorates \\
    to be defended publicly on \\
    \thesisdefencedate

    \bigskip
    \bigskip

    by

    \bigskip
    \bigskip

    {\headingfont\Large\bfseries \thesisauthorfirstnames\ \MakeUppercase{\thesisauthorlastname}}

    \bigskip
    \bigskip

    \thesisdegree \\
    born in \thesisbirthplace

    \vspace*{2\bigskipamount}
\end{center}
\makeatother

%% ---------------------------------------------------------------------
%% Page iv: doctoral committee and colophon
%% ---------------------------------------------------------------------
\clearpage
\thispagestyle{empty}
{\small
\setlength{\parskip}{0pt}
\noindent This dissertation has been approved by the promotors.

\medskip
\noindent Composition of the doctoral committee:

\smallskip
\noindent
\begin{tabular}{@{}p{0.5\textwidth}@{}p{0.5\textwidth}@{}}
    Rector Magnificus,  & chairperson \\
    \thesiscommittee
\end{tabular}

\ifx\thesiscommitteenote\empty\else
    \medskip
    \noindent \thesiscommitteenote
\fi

\vfill
\noindent\includegraphics[height=8mm, alt={Delft University of Technology logo}]{Images/_title_images/logo_black.pdf}

\bigskip
\ifx\thesisfunding\empty\else
    \noindent \thesisfunding

    \medskip
\fi
\ifx\thesisdataavailability\empty\else
    \noindent \thesisdataavailability

    \medskip
\fi
\noindent
\begin{tabular}{@{}p{0.22\textwidth}@{}p{0.78\textwidth}@{}}
    \textit{Keywords:}   & \thesiskeywords \\[\smallskipamount]
    \textit{Printed by:} & \thesisprinter  \\[\smallskipamount]
    \textit{Cover by:}   & \thesiscover    \\
\end{tabular}

\bigskip
\noindent Copyright \textcopyright\ \the\year\ by \thesisauthorinitials~\thesisauthorlastname
\ifx\thesislicence\empty\else\par\noindent \thesislicence\fi
\ifx\thesisorcid\empty\else

\smallskip
\noindent ORCID \href{https://orcid.org/\thesisorcid}{\thesisorcid}
\fi

\smallskip
\noindent ISBN \thesisisbn

\smallskip
\noindent An electronic copy of this dissertation is available at \\
\url{https://repository.tudelft.nl/}.
}
\cleardoublepage

\newpage

\begin{Spacing}{1.15} % Set up spacing for report

    %% =========== %%
    %% FRONTMATTER %%
    %% =========== %%
    \pagenumbering{roman}
    \pagestyle{frontmatter}
    \chapterstyle{frontmatter}
    \aliaspagestyle{chapter}{plain}

    \chapter*{Summary}
\addcontentsline{toc}{chapter}{Summary}

[A summary of the thesis in English: the problem, the approach, the main results and their significance. Typically one to three pages, without equations, figures or citations.]

\lipsum[1-2]

\chapter*{Samenvatting}
\addcontentsline{toc}{chapter}{Samenvatting}

\begin{otherlanguage}{dutch}
[Een samenvatting van het proefschrift in het Nederlands, met dezelfde inhoud als de Engelse samenvatting.]

\lipsum[1-2]
\end{otherlanguage}

    \chapter*{Preface}
\addcontentsline{toc}{chapter}{Preface}

% These instructions give you guidelines for preparing papers for AIAA Technical Papers using \LaTeX{}. Define all symbols used in the abstract. Do not cite references in the abstract. The footnote on the first page should list the Job Title and AIAA Member Grade for each author, if known Authors do not have to be AIAA members.

\lipsum[1-2]

    \cleardoublepage
    \setcounter{tocdepth}{1} % How deep the contents will actually go: 1 will only go into sections, 2 into subsection, 3 into subsubsections
    \tableofcontents*
    \cleardoublepage
    \listoffigures
    \cleardoublepage
    \listoftables
    \cleardoublepage
    % \newpage %% new page included due to the nomenclature package
    
    % Basically, reset counters every time you step into a new part of the report. Useful for when you have many many sections and you need to navigate quickly. Check the cls file for the numbering system.
    % \counterwithin*{equation}{section}
    % \counterwithin*{equation}{subsection}
    % \counterwithin*{equation}{subsubsection}
    % \counterwithin*{figure}{section}
    % \counterwithin*{figure}{subsection}
    % \counterwithin*{figure}{subsubsection}
    % \counterwithin*{table}{section}
    % \counterwithin*{table}{subsection}
    % \counterwithin*{table}{subsubsection}
    
    % Basically, creates new chapter numbering within a new part using the chngcntr package, this stops LaTeX from destroying the document outline when exporting, which is what happens when you \setcounter{chapter}{0} when the previous part had higher numbering.
    \counterwithin*{chapter}{part}
    
    % \input{Chapters/00Abstract.tex}
    
    %% Prints the nomenclature. Entries are defined in _nomenclature.tex.
\chapter*{Nomenclature}
\addcontentsline{toc}{chapter}{Nomenclature}
\markboth{Nomenclature}{Nomenclature}

%% List every defined entry. Remove this line to list only the entries used with \gls in the text.
\glsaddall

%% \printnom{list}{title}{style}{sort order}
%%   sort order: nocase (alphabetical by sort key) or def (order of definition)
\printnom{acronyms}{Acronyms}{nom-acronyms}{nocase}
\printnom{dimless}{Dimensionless numbers}{nom-definition}{nocase}
\printnom{latin}{Latin symbols}{nom-dimensions}{nocase}
\printnom{greek}{Greek symbols}{nom-dimensions}{def}
\printnom{super}{Superscripts}{nom-plain}{def}
\printnom{sub}{Subscripts}{nom-plain}{def}

% Finish the last page here, so the page numbering of the main matter starts on a fresh page.
\cleardoublepage

    
    %% ========== %%
    %% MAINMATTER %%
    %% ========== %%
    \pagenumbering{arabic}
    \chapterstyle{mainmatter}
    \aliaspagestyle{chapter}{empty}
    
    %% The abstract environment is used for the abstract of a chapter.
\newenvironment{chapterabstract}{%
    \list{}{\leftmargin\rightmargin}%
    \item%
    \relax%
    \rmfamily\itshape%
}{%
    \endlist%
}

\chapter{Introduction} \pagestyle{intro-biblio}
\chapterprecishere{
    \begin{epigraphs}
        \qitem{
            \textit{Nature and nature's laws lay hid in the night; God said `Let Newton be!' and all was light.}
        }{\opensans Alexander Pope} \\[1\baselineskip]
        \qitem{
            \textit{It did not last: the devil shouting `Ho. Let Einstein be!' restore the status quo.}
        }{\opensans Sir John Collings Squire}
    \end{epigraphs}
    
    \vspace{0pt plus 1filll} % fill space up to bottom

    \begin{chapterabstract}
        \lipsum[1]
    \end{chapterabstract}

    \blfootnote{Parts of this chapter have been published in Annalen der Physik \textbf{324}, 289 (1906)}
}
\chapterprecistoc{This command allows you to add a brief abstract to the ToC, in case you are interested in that sort of stuff. If not, then either leave this field empty, or don't bother calling upon this command.}

\lipsum[0-12]

 % actual pagestyle set in here

    \chapter{Examples} \pagestyle{mainmatter}
\chapterprecishere{}
\chapterprecistoc{}

\lipsum[1]

\section{Section}
\lipsum[1]

\subsection{Subsection}
\lipsum[1]

\subsubsection{Subsubsection}
\lipsum[1]

\paragraph{Paragraph}
\lipsum[1]

\subparagraph{Subparagraph}
\lipsum[1]

\raisebox{0pt}[\height][0pt]{\color{white} SEPARATOR TEXT - REMOVE WHEN DONE} \vspace{-2\baselineskip}



\section{Bibliography}

When there exists an author with last name prefix of 'van', use the \textbackslash VAN command, and follow it up with a \textasciitilde (see reference for example)

\citet*{Allen2002}

\citet*{Bettini2021}

\citet*{Eitelberg2020}

\citet*{Sinnige2020}

\citet*{Herrez2016}

\section{Tables}

\subsection{tabularx}

\begingroup
\renewcommand\tabularxcolumn[1]{m{#1}} % https://latex.org/forum/viewtopic.php?t=32433

\begin{minipage}[t]{0.48\textwidth}
\begin{table}[H]
    \centering
    \caption{Fuselage geometry \cite[Tab.~4]{Bettini2021}}
    \label{tab:fuselage_geom}
    {\footnotesize
    \begin{tabularx}{\textwidth}[t]{X|rl} \toprule
        \multicolumn{1}{c|}{\textbf{Parameter}} &
        \multicolumn{2}{c}{\textbf{Value}}                  \\ \midrule
        Length          & 1.342     & \si{\meter}           \\
        Diameter        & 0.140     & \si{\meter}           \\
        Volume          & 0.0160632 & \si{\meter\cubed}     \\
        Volume of aft support strut (estimate)  & 0.0004491    & \si{\meter\cubed}  \\
        Volume of wing support struts (estimate for both struts together)                                & 0.0035296   & \si{\meter\cubed} \\ \bottomrule
    \end{tabularx}
    }
\end{table}
\end{minipage}
\hfill
\begin{minipage}[t]{0.48\textwidth}
\begin{table}[H]
    \centering
    \caption{Propeller geometry \cite[Tab.~8]{Bettini2021}}
    \label{tab:prop_geom}
    {\footnotesize
    \begin{tabularx}{\textwidth}[t]{X|rl} \toprule
        \multicolumn{1}{c|}{\textbf{Parameter}}         &
        \multicolumn{2}{c}{\textbf{Value}}                              \\ \midrule
        Number of blades            & 6                 &               \\
        Diameter                    & 0.2032            & \si{\meter}   \\
        Pitch angle at $r/R = 0.7$  & 45                & deg           \\
        Rotation direction (from rear)  & CW    & 2x \\
                                        & CCW   & 1x \\ \bottomrule 
    \end{tabularx}
    }
\end{table}
\end{minipage}

\begin{table}[H]
\caption{Wing geometry \cite[Tab.~5]{Bettini2021}}
\label{tab:wing_geom}
    \begin{subtable}[t]{0.48\textwidth}
    \centering
    {\footnotesize
    \begin{tabularx}{\textwidth}[t]{X|rl} \toprule
        \multicolumn{1}{c|}{\textbf{Parameter}}         &
        \multicolumn{2}{c}{\textbf{Value}}                          \\ \midrule
        Span                    & 1.397     & \si{\meter}           \\
        Area                    & 0.2172    & \si{\meter\squared}   \\
        Mean aerodynamic chord  & 0.165     & \si{\meter}           \\
        Aspect ratio            & 8.98      &                       \\
        Taper ratio             & 0.40      &                       \\
        Sweep angle at $0.25c$  & 0         & deg                   \\
        Incidence angle         & 0         & deg                   \\
        Dihedral angle          & 4         & deg                   \\
        Twist                   & 2         & deg                   \\ \bottomrule
    \end{tabularx}
    }
    \end{subtable}
    \hfill
    \begin{subtable}[t]{0.48\textwidth}
    \centering
    {\footnotesize
    \begin{tabularx}{\textwidth}[t]{X|rl} \toprule
        \multicolumn{1}{c|}{\textbf{Parameter}}         &
        \multicolumn{2}{c}{\textbf{Value}}                                                  \\ \midrule
        Airfoil (constant over span)        & DU 96-150                 &                   \\
        Aileron span                        & 0.197                     & \si{\meter}       \\
        Aileron spanwise position           & 0.71 < y/(b/2) < 0.99     &                   \\
        Aileron chord                       & 0.029                     & \si{\meter}       \\ 
        Flap span (per side)                & 0.369                     & \si{\meter}       \\
        Flap chord                          & 0.060                     & \si{\meter}       \\
        Wing $0.25c$ from fuselage nose     & 0.680                     & \si{\meter}       \\
        Volume                              & 0.0030229                 & \si{\meter\cubed} \\ \bottomrule
    \end{tabularx}
    }
    \end{subtable}
\end{table}

\endgroup


\section{Graphs}

\subsection{subfigures}
\begin{figure}[H]
    \centering
    \begin{subfigure}[t]{0.48\textwidth}
        \centering
        \includegraphics[width=1.0\textwidth]{Images/graph.jpg}
        \caption{graph 1}
        \label{subfig:graph1}
    \end{subfigure}
    \hfill
    \begin{subfigure}[t]{0.48\textwidth}
        \centering
        \includegraphics[width=1.0\textwidth]{Images/graph.jpg}
        \caption{graph 2}
        \label{subfig:graph2}
    \end{subfigure}
    \bigskip
    \begin{subfigure}[t]{0.48\textwidth}
        \centering
        \includegraphics[width=1.0\textwidth]{Images/graph.jpg}
        \caption{graph 3 help help help help help help help help help help help help help}
        \label{subfig:graph3}
    \end{subfigure}
    \hfill
    \begin{subfigure}[t]{0.48\textwidth}
        \centering
        \includegraphics[width=1.0\textwidth]{Images/graph.jpg}
        \caption{graph 4}
        \label{subfig:graph4}
    \end{subfigure}
    \caption{graphs (subfigures)}
    \label{fig:graph_subfigures}
\end{figure}

\subsection{Continued float}
\begin{figure}[H]
    \centering
    \begin{subfigure}[t]{0.48\textwidth}
        \centering
        \includegraphics[width=1.0\textwidth]{Images/graph.jpg}
        \caption{graph 5}
        \label{subfig:graph5}
    \end{subfigure}
    \hfill
    \begin{subfigure}[t]{0.48\textwidth}
        \centering
        \includegraphics[width=1.0\textwidth]{Images/graph.jpg}
        \caption{graph 6}
        \label{subfig:graph6}
    \end{subfigure}
\end{figure}%
\begin{figure}[H]\ContinuedFloat
    \begin{subfigure}[t]{0.48\textwidth}
        \centering
        \includegraphics[width=1.0\textwidth]{Images/graph.jpg}
        \caption{graph 7}
        \label{subfig:graph7}
    \end{subfigure}
    \hfill
    \begin{subfigure}[t]{0.48\textwidth}
        \centering
        \includegraphics[width=1.0\textwidth]{Images/graph.jpg}
        \caption{graph 8}
        \label{subfig:graph8}
    \end{subfigure}
    \caption{graphs (continued float)}
    \label{fig:graph_contfloat}
\end{figure}

\subsubsection{Aligned graphs}

\begin{figure}[H]
\begin{tabular}{@{}p{0.485\textwidth}p{0.485\textwidth}@{}}
    \raisebox{-\height}{\includegraphics[width=1.0\linewidth]{Images/graph.jpg}}&
    \raisebox{-\height}{\includegraphics[width=1.0\linewidth]{Images/graph.jpg}}\\
    \caption{graph}\label{fig:graph9}&
    \caption{graph graph graph graph graph graph graph graph graph graph graph graph graph graph graph graph}\label{fig:graph10}
\end{tabular}
\end{figure}

\subsubsection{Aligned repeated graphs}
\begin{figure}[H]
\begin{tabular}{@{}p{0.485\textwidth}p{0.485\textwidth}@{}}
    \raisebox{-\height}{\includegraphics[width=1.0\linewidth]{Images/graph.jpg}}&
    \raisebox{-\height}{\includegraphics[width=1.0\linewidth]{Images/graph.jpg}}\\
    \repeatcaption{fig:graph9}{graph}&
    \repeatcaption{fig:graph10}{graph graph graph graph graph graph graph graph graph graph graph graph graph graph graph graph }
\end{tabular}
\end{figure}

\begin{landscape}
\section{Landscape environment}

The landscape environment basically rotates your text 90 degrees counterclockwise, while also rotating the Overleaf environment so you can actually see what you are writing. You can insert a sequence of really long equations in here:

\begin{alignat*}{3}
    a &\,\,=\,\,&& \pderiv[]{a}{x} \pderiv[]{x}{y} \pderiv[]{y}{z} \pderiv[]{y}{z} \pderiv[]{y}{z} \pderiv[]{y}{z} \pderiv[]{y}{z} \pderiv[]{y}{z} \pderiv[]{y}{z} \pderiv[]{y}{z} \pderiv[]{y}{z} \pderiv[]{y}{z} \pderiv[]{y}{z} \pderiv[]{y}{z} \pderiv[]{y}{z} \pderiv[]{y}{z} \pderiv[]{y}{z} \pderiv[]{y}{z} \pderiv[]{y}{z} + \pderiv[]{y}{z} \pderiv[]{y}{z} \pderiv[]{y}{z} \pderiv[]{y}{z} \pderiv[]{y}{z} \pderiv[]{y}{z} \pderiv[]{y}{z} \pderiv[]{y}{z} \pderiv[]{y}{z} \pderiv[]{y}{z} \pderiv[]{y}{z} \pderiv[]{y}{z} \\[3mm]
    b &\,\,=\,\,&& a \exp{\big\{a + b + c + d + e + f + g + h + i\big\}}
\end{alignat*}

Alternatively, you can also insert a table that may be really wide like this:

\begin{tabularx}{\linewidth}{r|X}
    \bfseries first column & \bfseries a really really really really really really really really really long column \\ \hline
    1            & 2    \\ \hline
    1            & 2    \\ \hline
    1            & 2    \\ \hline
    1            & 2    \\ \hline
    1            & 2 
\end{tabularx}

\clearpage

You can also see here that this will rotate the text of the opposing page (recto if previous page was verso, and verso if previous page was recto) in the same orientation so you only need to rotate the book once to read the text/graphs/tables/equations: you do not need to flip the book/report.

\clearpage

and this is what it looks like when the last landscape scenario ensures (either a verso-recto pair, or a recto-verso pair)
\end{landscape}

\section{tmp} % actual pagestyle set in here

    \chapterstyle{bibliography}
    \aliaspagestyle{chapter}{plain}
    \DeclareRobustCommand{\VAN}[0]{van} % Change from Van to van
    % \printbibliography[title={References}, heading=bibintoc]
    \renewcommand{\bibname}{References}
    %\begin{multicols}{2}[\printbibheading]
    \setlength\bibitemsep{0.6\baselineskip} %% half of Spacing (so if 1.2 spacing -> 0.6\baselineskip)
    \printbibliography%[heading=none]
    \pagestyle{intro-biblio}
    %\end{multicols}
    \DeclareRobustCommand{\VAN}[0]{Van} % Change from van to Van for appendix purposes

    %% ======== %%
    %% APPENDIX %%
    %% ======== %%
    \bookmarksetup{startatroot} % https://tex.stackexchange.com/questions/131959/pdf-bookmarks-for-references-and-index-appear-as-contained-in-the-appendix-wron/131966
    \chapterstyle{appendix}
    \aliaspagestyle{chapter}{empty}
    \appendix

    \clearpage
    \chapter{TODO} \label{app:litStudy} \pagestyle{appendix}
\chapterprecishere{}
\chapterprecistoc{}

TODO


    
    \clearpage
\end{Spacing}

\end{document}